\documentclass[journal]{IEEEtran}
\usepackage[margin=1in]{geometry}
\pdfinfoomitdate=1
\pdftrailerid{}
\pdfsuppressptexinfo=-1
\usepackage{amsmath,amssymb,amsthm,bm,booktabs,graphicx}
\usepackage{xr-hyper}
\usepackage[hidelinks]{hyperref}
\externaldocument{main}[main.pdf]
\newcommand{\E}{\mathbb E}
\newcommand{\tr}{\operatorname{tr}}
\newcommand{\Pp}{\mathbb P}
\newcommand{\KL}{\operatorname{KL}}
\newtheorem{proposition}{Proposition}
\title{Supplementary Material: Learning-Aware Lyapunov Resource Control for Correlated Multi-Agent Markov Learning}
\author{Jialin Zhuang and Vincent K. N. Lau}
\begin{document}
\maketitle
\renewcommand{\theequation}{S\arabic{equation}}

\section{Proof of the Same-Stream Actual-Risk Enclosure}
Consider the stationary common factor in the main paper, with independent private observation noise and a deterministic probe $d=(q_d,b_d,n_d)$ ending at $T=n_db_d$.
The probe sigma-field $\mathcal H_d$ contains its observations; spatial contrasts contain no additional information about the common factor.
Under true instance $j$, Gaussian filtering gives
\begin{align}
C_T\mid\mathcal H_d&\sim\mathcal N(m_j,P_j),\nonumber\\
m_j&\sim\mathcal N(0,s_j^2),\quad s_j^2=\theta_j-P_j.
\end{align}
The posterior variance $P_j$ is deterministic because the probe schedule is fixed before data.
Starting from variance $\theta_j$, the observation update is $P\leftarrow P/(1+q_dP)$; between observations it is $P\leftarrow\lambda^{2b_d}P+(1-\lambda^{2b_d})\theta_j$.
The last update has no subsequent prediction.

Fix either downstream program $a$, including its budget, initialization, sample spacing and deterministic weights.
All learning observations occur strictly after $T$.
With $r_a=\lambda^{b_a}$ and the notation of main Eq.~\eqref{eq:linear-program}, its unconditional stationary risk is
\begin{align}
J_{j,a}
&=d_a^2+\theta_j\sum_{t,u=1}^{N_a}\omega_{a,t}\omega_{a,u}r_a^{|t-u|}\nonumber\\
&\quad+q_a^{-1}\sum_{t=1}^{N_a}\omega_{a,t}^2.
\label{eq:supp-exact-risk}
\end{align}
Conditional means of future common states are $\lambda^{b_at}m_j$.
Their conditional covariance is their stationary covariance minus $s_j^2\lambda^{b_a(t+u)}$.
Private learning noise is independent of the probe.
Therefore
\begin{equation}
\E_j[Z_a^2\mid\mathcal H_d]
=J_{j,a}+2d_ak_am_j+k_a^2(m_j^2-s_j^2).
\label{eq:supp-conditional-risk}
\end{equation}
This identity includes the posterior mean, which is unbounded across probe histories; no uniform conditional Gaussian bound is invoked.

Let $E_j$ be the event that the decision selects the wrong frozen program, and let $p_j=\Pr_j(E_j)\le\delta$.
Subtracting Eq.~\eqref{eq:supp-conditional-risk} for the correct and wrong programs gives
\begin{equation}
\E_jZ_{\widehat a}^2=J_{j,a_c}+p_j\Delta J_j+\E_j[\mathbf1_{E_j}U_j],
\end{equation}
where
\begin{align}
U_j&=2(d_{a_w}k_{a_w}-d_{a_c}k_{a_c})m_j\nonumber\\
&\quad+(k_{a_w}^2-k_{a_c}^2)(m_j^2-s_j^2).
\end{align}
The centered Gaussian first and second moments are orthogonal, with variances $s_j^2$ and $2s_j^4$.
Consequently $\E_jU_j=0$ and $\E_jU_j^2=\Gamma_j^2$ as defined in the main paper.
Cauchy--Schwarz gives $|\E_j[\mathbf1_{E_j}U_j]|\le\sqrt{p_j}\Gamma_j\le\sqrt\delta\Gamma_j$.
Bounding the signed term $p_j\Delta J_j$ and using nonnegative squared risk proves main Theorem~\ref{thm:same-stream}.

\section{Exact Weights of the Implemented Learners}
Let $M_a=\min\{\lfloor B_m'/(h+q_a)\rfloor,\lfloor B_e'/b_a\rfloor\}$ be the number of scheduled observations after subtracting probe costs.
All $M_a$ observations are paid.
The FIFO withholding convention leaves $N_a=(M_a-D)_+$ usable observations, not a free replacement budget.
For $N_a=0$ the reported parameter is one, so $d_a=1$ and the weight vector is empty.

\subsection{Running mean}
The Layer-A output is $(1+\sum_{t=1}^{N_a}\bar X_t)/(N_a+1)$.
Thus $d_a=1/(N_a+1)$ and every $\omega_{a,t}=1/(N_a+1)$.
Main Eq.~\eqref{eq:linear-program} then gives
\begin{equation}
k_a=\frac{r_a(1-r_a^{N_a})}{(N_a+1)(1-r_a)}.
\end{equation}

\subsection{Averaged FIFO affine learner}
In Layer B, an increment $\eta(\bar X_t-\mu w_{t-1})$ is computed before enqueueing and is applied after $D$ opportunities.
The original implementation fixes $\mu=1/2$ and $\eta=0.1/[\mu(D+1)]$, starts from parameter one with an empty queue, and averages the applied parameters.
Indexing these applied parameters as $z_n$ gives
\begin{equation}
z_n=z_{n-1}-\eta\mu z_{n-D-1}+\eta\bar X_n,
\quad z_{-D}=\cdots=z_0=1.
\end{equation}
Define $h_0=1$, $h_t=0$ for $t<0$, and $h_t=h_{t-1}-\eta\mu h_{t-D-1}$ for $t\ge1$.
For $N=N_a>0$, the output $N^{-1}\sum_{n=1}^Nz_n$ has exact coefficients
\begin{equation}
\omega_{a,t}=\frac\eta N\sum_{v=0}^{N-t}h_v,
\quad d_a=1-\mu\sum_{t=1}^N\omega_{a,t}.
\label{eq:supp-fifo-weights}
\end{equation}
To verify the deterministic term, constant input $\bar X_n=\mu$ keeps the initial history at one; linearity then yields the identity for $d_a$.
Eq.~\eqref{eq:supp-fifo-weights} represents the average of applied iterates; all scheduled observations remain charged.

\section{Shrinking-Risk and Budget-Shift Interface}
Suppose the finite action catalogue has $|\lambda|\le1-\gamma<1$, bounded $D$, and strictly stable affine recursions.
For the registered step rule this stability follows directly from the characteristic polynomial $z^{D+1}-z^D+t$, $t=\eta\mu$.
At small $t>0$, the root at one moves inside the unit disk and the other roots remain inside.
Writing a first unit-circle crossing as $z=e^{i\omega}$ yields $t=2\sin(\omega/2)$ and the smallest positive crossing $\omega=\pi/(2D+1)$.
Thus $0<t<2\sin[\pi/(4D+2)]$ is stable; the registered $t=0.1/(D+1)$ lies strictly inside this interval.
The impulse response consequently decays geometrically, and $\sum_{v\ge0}h_v=1/(\eta\mu)$.

The interior weights in Eq.~\eqref{eq:supp-fifo-weights} are $1/(\mu N)$ up to a summable end correction: their difference from these constant weights has $\ell_1$ norm $O(N^{-1})$ and maximum absolute value $O(N^{-1})$.
The homogeneous transient is summable, so its average $d_a$ is $O(N^{-1})$.
The AR(1) covariance kernel has uniformly bounded absolute row sums.
Substitution into Eq.~\eqref{eq:supp-exact-risk} therefore gives
\begin{equation}
J_{j,a}=\frac1{\mu^2N_a}\left[\frac1{q_a}+\theta_j\frac{1+r_a}{1-r_a}\right]+O(N_a^{-2}).
\label{eq:supp-averaged-rate}
\end{equation}
The running mean has the same expansion with $\mu=1$.
These constants are uniform over the finite catalogue and bounded instance class.
On a budget ray, $N_a=s/c_a+O(v+D+1)$ with $c_a=\max\{(h+q_a)/\beta_m,b_a/\beta_e\}$ and fixed probe charge $v$.
Hence $J_{j,a}=K_{j,a}/s+O((v+D+1)/s^2)$, with
\begin{equation}
K_{j,a}=\frac{c_a}{\mu^2}\left[\frac1{q_a}+\theta_j\frac{1+r_a}{1-r_a}\right].
\end{equation}
Subtracting a fixed probe charge changes $N_a$ by at most $O(v+D+1)$, giving the claimed budget-shift remainder.
Moreover $k_a=O(N_a^{-1})$ because the weights are uniformly $O(N_a^{-1})$ and $\sum_{t\ge1}|r_a|^t$ is bounded.
Since $s_j^2\le\theta_j$, main Eq.~\eqref{eq:selection-carryover} yields
\begin{equation}
\sqrt\delta\Gamma_j=O(\sqrt\delta/s^2).
\end{equation}
For a pair of frozen actions satisfying the main phase theorem's leading-gap and wrong-decision assumptions, adding this remainder to admission changes only its second-order class constants.
The same necessary-information and sufficient-scale argument then applies to the corrected actual-risk rule.
This is an averaged, additive Gaussian affine result; it does not remove the nonzero noise floor of a fixed-step unaveraged learner or the conditional-history requirements of the general Markov theorem.

\section{Complete Existing-Evidence Reanalysis}
The calculation uses all registered EXP-016B budget points, both regimes and all eight policies: 14,336 policy cells, each with its original 192-seed empirical mean.
Layer A contains 104 scenario families and Layer B its original 56, including eight search-censored families in each layer.
For fixed policies Eq.~\eqref{eq:supp-exact-risk} is exact.
For probe policies main Eq.~\eqref{eq:same-stream-risk} supplies population-risk bounds.
When information-only and learning-aware policies execute the same frozen program, their paired difference is identically zero, not the difference of two independent interval bounds.
Otherwise subtract their one-sided bounds.

\begin{table}[t]
\centering
\caption{Population enclosures for mean information-only minus learning-aware risk, rounded outward.
All means use the listed complete stratum.}
\begin{tabular}{llrrr}
\toprule
Layer & Stratum & Cells & Lower & Upper\\
\midrule
A & All & 1184 & .0506 & .0706\\
A & Intermediate & 576 & .0770 & .1180\\
A & Original practical & 446 & .1607 & .1952\\
A & Search-censored & 32 & .4879 & .4880\\
B & All & 608 & .0079 & .0223\\
B & Intermediate & 288 & -.0049 & .0252\\
B & Original practical & 222 & .0562 & .0773\\
B & Search-censored & 32 & .1957 & .1958\\
\bottomrule
\end{tabular}
\end{table}
Below identification, all 208 Layer-A and 112 Layer-B comparisons are identical programs; above finite value thresholds the corresponding counts are 384 and 192.
In the original practical subset, 323/446 Layer-A and 176/222 Layer-B cells have strictly positive improvement lower bounds, while 101 and 33 have strictly negative upper bounds; the remaining 22 and 13 are inconclusive.
These include regime-specific losses even when the stratum-average improvement is positive.
No new admission rule was executed and no old gate was replaced.
The old probe-length routine used the sufficient exponent $\log(0.5/0.025)$, giving an equal-prior average-error target of $0.025$ rather than a directional guarantee of $0.025$.
The reanalysis uses each probe's exact affinity $\exp(-I_B)$ as a valid bound on either directional error; over all 104 frozen probe designs these bounds range from $0.0352271$ to $0.0497465$.
Failure of this sufficient directional certificate does not establish that the actual directional errors exceed $0.025$.
No outcomes or probe lengths are changed to obtain a smaller error bound.
This outcome-known reanalysis gives population enclosures, not new confidence intervals for empirical means.

\section{Complete Threshold-Sandwich Scale Calculation}
For the independently initialized Gaussian running mean, let $c_a=\max\{(h+q)/\beta_m,b/\beta_e\}$ and $r=\lambda^b$.
The exact horizon is $N_a=s/c_a+O(1)$, with bounded terminal charges and rounding.
Main Eq.~\eqref{eq:mean-risk} and Eq.~\eqref{eq:geometric-correlation-sum} give $K_j(a)=c_a\{q^{-1}+\theta_j(1+r)/(1-r)\}$ and a uniform $O(s^{-2})$ remainder.
Removing probe charge $v$ changes $N_a$ by $O(v+1)$ and hence risk by $O((v+1)/s^2)$ for $s\ge s_0(v+1)$; compact persistence and the finite catalogue make these constants uniform.
Thus the phase-theorem risk conditions have a concrete Gaussian realization, separately from the constant-step affine certificate.
Write $d_\delta=\operatorname{kl}(1-\delta,\delta)$, $r_0=1-\gamma<1$, and $Q=\max{\cal Q}$.
Under either true instance, predictable spacing gives $\|C_{S_k}\|_{L^2}\le r_0\|C_{S_{k-1}}\|_{L^2}+\sqrt{\theta_+}$, hence a uniform bound $M_C=\sqrt{\theta_+}/(1-r_0)$.
The normalized observation $Y_k/\sqrt{q_k}$ has $L^2$ norm at most $M_C+1$.
For either candidate filter, write its posterior mean as $(1-t_{j,k})m^-_{j,k}+t_{j,k}Y_k/\sqrt{q_k}$, with predictable $0\le t_{j,k}\le1$.
Prediction multiplies this mean by a factor of magnitude at most $r_0$; the triangle inequality bounds both predictive means uniformly in $L^2$ by $M_F=r_0(M_C+1)/(1-r_0)$, starting from zero.
Both predictive variances are at most $\theta_+$, so $1\le V_{j,k}\le V_+=1+Q\theta_+$.
Main Eq.~\eqref{eq:conditional-kl} gives the expected, rather than pointwise, bound $\E_j\mathcal I_{j,1-j}(Z_k,A_k)\le i_+$, where $i_+=(\log V_++V_++4QM_F^2)/2$ is finite.
Extend a stopped rule predictably after stopping; nonnegativity and its deterministic budget bound $n_{\max}(s)\le c_Bs$ imply $\E_j\sum_{k\le\tau}\mathcal I_{j,1-j}\le i_+n_{\max}(s)$.
Together with main Eq.~\eqref{eq:info-constraint}, this proves $B_N\ge c_Nd_\delta$.
The fixed probe has information at least $i_->0$, so $n_\delta=\lceil i_-^{-1}\log(1/\delta)\rceil$ observations suffice and cost at most $c_pn_\delta+D_{\max}$.
Continuity of $\log(1/\delta)/d_\delta$ gives $n_\delta\le c_0+c_1d_\delta\le c_0+c_2B_N$.
Write $u_\delta\le c_pn_\delta+D_{\max}$ for the full probe charge.
The assumed expansion and budget-shift inequality give a high-instance improvement over the full-budget fallback certificate of at least $g_{\min}/s-\{2M+L(u_\delta+1)\}/s^2-\delta w_+/s$.
It is strictly positive for $s>\{4M+2L(u_\delta+1)\}/g_{\min}$, also requiring $s\ge s_0(u_\delta+1)$.
On instance zero, the assumed coefficient ordering gives $\epsilon_{\rm safe}\le\{2M+L(u_\delta+1)\}/s^2+\delta_0w_+/s$.
Thus $s\ge\sqrt{2\{2M+L(u_\delta+1)\}/\bar\epsilon}+2\delta_0w_+/\bar\epsilon$ suffices.
Using $2\sqrt{xy}\le x+y$, all three requirements are bounded by a class-constant multiple of $n_\delta+B_{\rm oracle}+B_\epsilon+1$.
Taking the maximum proves the upper bound in main Eq.~\eqref{eq:threshold-sandwich}; its lower bound follows because every admitted $\delta$-correct probe obeys both information constraints.

\section{Robust-Certificate Sensitivity Details}
This section supplies the proof of main Corollary~\ref{cor:robust-certificate}.
Write $\delta^+=C_q(r_q^+)^b$.
On the coverage event, $\delta^+\ge\delta$ and monotonicity on $\mu-2L\delta^+>0$ gives $a^+\ge a$ and $\beta^+\ge\beta$.
Evaluating the delayed Young inequality at the shared $\lambda^+=\sqrt{h/a^+}$ proves the upper recursion; separately optimized forcing terms need not be monotone.
For $\Delta\delta=\delta^+-\delta$, $\Delta K=K_q^+-K_q$, and $\Delta\Omega=\Omega_q^+-\Omega_q$,
\begin{align}
\Delta a
&=2\eta^2\Delta K+(2\eta L+4\eta^2L^2)\Delta\delta,\nonumber\\
\Delta\beta
&=2\eta^2\Delta\Omega+4\eta^2G^2\Delta\delta\nonumber\\
&\quad+4\eta G^2\left[
\frac{(\delta^+)^2}{\mu-2L\delta^+}
-\frac{\delta^2}{\mu-2L\delta}
\right],\nonumber\\
\Delta c
&=(\sqrt{a^+}+\sqrt h)^2-(\sqrt a+\sqrt h)^2.
\end{align}
For $h>0$,
\begin{align}
\Delta d
&\le(1+\sqrt{h/a^+})\Delta\beta
+\sqrt h\left(\frac1{\sqrt a}-\frac1{\sqrt{a^+}}\right)\beta\nonumber\\
&\quad+\frac{\sqrt{a^+}-\sqrt a}{\sqrt h}g,
\end{align}
and $\Delta d=\Delta\beta$ when $h=0$.
For $U(c,d)=d/(1-c)+c^n(R_0-d/(1-c))_+$, $|\partial_dU|\le(1-\bar c)^{-1}$ and $|\partial_cU|\le nR_0+\bar d(1-\bar c)^{-2}$.
The mean-value bound proves main Eq.~\eqref{eq:robust-score-gap}; applying it to the upper- and exact-score minimizers gives $U_{\widehat a}\le U_{a^\star}+2\max_a\varepsilon_a$.

\section{Stochastic-Delay Interpretation}
The controlled stochastic delay differential equation associated with one fixed resource action is
\begin{equation}
dE(t)=-A_qE(t-\tau(t))dt
+\sqrt{\eta}\,\Sigma_q^{1/2}(E_t)dW(t).
\end{equation}
Here $\Sigma_q(E_t)$ retains both additive innovation covariance and the state-dependent quadratic variation of random temporal-difference Jacobians.
For $V(E)=\|E\|^2$, It\^o's formula gives
\begin{equation}
{\cal L}_aV(E_t)
=-2E(t)^\top A_qE(t-\tau(t))
+\eta\tr\Sigma_q(E_t).
\end{equation}
Participation controls spatial covariance, spacing controls mixing, and $\eta$ controls their scale; the discrete theorem retains finite steps, integer delays and rounding.
\section{Constructing a Certificate from a Known Finite Model}
Let $Z$ be the finite augmented state containing all sources needed for the retained joint TD tuple.
For each participation $q$, suppose its kernel $M_q$ and stationary law $\pi_q$ are known, and the marginal target policy and reward law are preserved.
Tabulate $H_i(z)$ and $c_i(z)$, put $A=\E_{\pi_q}H_i$, $c=\E_{\pi_q}c_i$, and solve $Aw^\star=c$.
Check that these means are common to all actors and all $q$, and that $\mu=\lambda_{\min}((A+A^\top)/2)>0$.
With $\xi_i(z)=c_i(z)-H_i(z)w^\star$, compute
\begin{align}
L&=\max_{i,z}\|H_i(z)\|_2,\qquad G=\max_{i,z}\|\xi_i(z)\|,\nonumber\\
K_q&=\lambda_{\max}\!\left(\sum_z\pi_q(z)\bar H_q(z)^\top\bar H_q(z)\right),\nonumber\\
\Omega_q&=\sum_z\pi_q(z)\|\bar\xi_q(z)\|^2,\nonumber\\
\beta_q(g)&=\max_z\frac12\sum_{z'}|M_q^g(z,z')-\pi_q(z')|.
\label{eq:finite-model-table}
\end{align}
Here bars denote actor averages, not averages of squared norms.
The integer $g$ is the number of transitions of the augmented tuple chain between the conditioning state and the retained tuple; it must be derived from the declared physical schedule, not identified with a state lag or drain charge.

\begin{proposition}[Constructive finite-model interface]
If $M_q$ is irreducible and aperiodic and the above common-mean and positivity checks pass, Assumption~\ref{ass:catalogue}'s boundedness, centering and conditional mixing inputs hold with Eq.~\eqref{eq:finite-model-table} and $\delta=\beta_q(g)$.
For every action whose resulting $\mu_\delta$, $a_\delta$ and $c_{\rm aff}<1$ pass the main paper's shield, the affine certificate is valid from any declared initial augmented state.
\end{proposition}
\begin{proof}
The finite maxima bound every sample and stationarity gives $\E_{\pi_q}\xi_i=0$.
Conditional on any earlier history, the Markov property expresses the next tuple law as a mixture of rows of $M_q^g$.
Convexity of total variation bounds its distance from $\pi_q$ by $\beta_q(g)$, uniformly over that history.
The drift proof uses only this total-variation distance, so the geometric envelope can be replaced by this exact finite-gap value.
The stated positivity and contraction checks are precisely the remaining hypotheses of the affine bound.
\end{proof}

This is a construction, not a statistical estimator of an unknown kernel.
A dense joint-state table may grow exponentially with participation; its offline construction cost is separate from the online scalar score.
In the seven-state shared/private-source model, a product coupling gives the cheaper bound
$\beta_q(g)\le\min\{1,(q+1)\beta_{\rm pair}(g)\}$, where
$\beta_{\rm pair}(g)=\max_x\|M_{\rm pair}^g(x,\cdot)-\pi_{\rm pair}\|_{\rm TV}$.
This uses the actual consecutive-pair chain and independent latent source chains; fresh selector coins are integrated by data processing.
The existing implementation computes these pair kernels and the same-time aggregate moments directly.
For a moving target policy, a frozen-kernel table does not certify mixed histories: either supply uniform bounds for the nonhomogeneous kernel products or reset the learner and exclude all pre-reset target updates.

\section{Executable Actual-Risk Gate and Saved-Result Reconstruction}
For each fixed candidate probe, Gaussian filtering gives $P_j$, the formulas above give $(J_{j,a},d_a,k_a)$, and the probe likelihood laws give a directional error bound $\delta_j$ under each hypothesis.
The exact Gaussian affinity is one valid choice: a zero-threshold likelihood-ratio decision obeys $\Pr_j(\mathrm{error})\le\int\sqrt{p_0p_1}$ in each direction.
An equal-prior average-error bound is not interchangeable with this directional bound.
Compute both hypothetical-instance quantities before probing; the actual unknown instance is not a gate input.

\begin{proposition}[Actual-risk admission]
Fix the two hypotheses, noise law, learner initialization and deterministic linear program weights as in Theorem~\ref{thm:same-stream}.
Given a finite public probe catalogue, choose a feasible design satisfying Eq.~\eqref{eq:actual-admission} before any probe observations, with deterministic tie breaking.
After its likelihood decision, execute its frozen selected program on the continuing stream.
Then $\E_1Z^2<F_1$ and $\E_0Z^2\le F_0+\bar\epsilon$.
If none passes, no probe is executed and the actual risk is exactly $F_j$.
\end{proposition}
\begin{proof}
For each fixed design Theorem~\ref{thm:same-stream} gives $\E_jZ^2\le U_j(d)$.
Pre-observation deterministic choice leaves its fixed-design hypothesis intact.
The two gate inequalities imply the conclusions. In the rejection branch the stationary full-budget commitment defines $F_j$.
\end{proof}
Thus the procedure is: price the probe, compute the remaining horizons and exact weights for both branch programs, compute $U_0,U_1,F_0,F_1$, apply the two inequalities, and only then acquire data.
An observation-adaptive choice of probe duration or of program weights is not covered by this fixed-design argument.

We reconstruct this gate on the existing EXP-016B outcomes, using its single registered information-only probe candidate and the original scenario-specific allowances $0.1$ or $0.2$.
All $896$ public layer/scenario/budget designs and both regimes are retained, including search-censored families.
The corrected gate admits $480$ designs, versus $505$ without the carryover term; it removes $25$ admissions.
The outcome-known zero-allowance stress case admits none under either rule.
It is retained separately: a paid probe can reduce the low-instance learning horizon even with perfect identification.

\begin{table}[t]
\centering
\caption{Saved-result gate reconstruction: mean squared risk over all budget/regime cells, without new trajectories. Each original cell mean contains all 192 saved seeds.}
\label{tab:gate-reconstruction}
\begin{tabular}{lrrrr}
\toprule
& A high & A low & B high & B low\\
\midrule
Full-budget fallback & .63455 & .09201 & .82703 & .54369\\
Information-only & .54814 & .24062 & .76422 & .62515\\
Original learning-aware & .57637 & .09304 & .79121 & .55825\\
No-carryover gate & .52548 & .09488 & .74204 & .56252\\
Corrected gate & .52989 & .09393 & .74242 & .56068\\
\bottomrule
\end{tabular}
\end{table}

The gate is a deterministic function of public design quantities, so an admitted design takes the saved information-only outcome and a rejected design takes the saved no-probe fallback outcome.
The simulator's potential observations are policy independent and shared across these branches; this reconstructs their empirical means without generating or resampling a trajectory.
Table~\ref{tab:gate-reconstruction} is therefore a retrospective policy comparison, not fresh execution or independent confirmation.
Against fallback, the corrected gate retains $13/592$ and $4/304$ empirically adverse high-instance cells in A and B; finite-sample cell means need not satisfy the population ordering.
Its low-instance mean risks exceed fallback by $.00192$ and $.01699$, respectively, within the declared nonzero safety allowances, not a claim of zero-harm safety.
The no-carryover ablation has slightly lower high-instance means but higher low-instance means; the correction trades admissions for a valid population-risk guarantee, not uniformly better observed risk.
All eight original policies, per-cell costs, both replay variants and their complete cell tables are preserved.
The new comparison was specified after earlier outcomes were known and supplies no new confidence interval.

\section{Nearest-Work Assumptions and Objectives}
The comparison uses the accessible full preprints, with theorem and proof locators recorded in the accompanying source-reading matrix; it is not an exhaustive priority claim.
DASA~\cite{dalfabbro2024dasa}, Assumption~3 and Theorem~1, combines independent agent chains with median-based update filtering to obtain average-delay dependence and parallel speedup.
Its Lemma~2 uses cross-agent independence in the noise cross terms.
Our bounded deterministic delay condition is narrower; our added object is covariance-dependent paid learning risk, not a stronger delay rate.

Perazzone et al.~\cite{perazzone2025lyapunov}, Theorem~1 and Eqs.~(12)--(23), link inverse participation probabilities to a nonconvex convergence bound, then optimize a time-average convergence/communication surrogate with power queues.
Their solver and instantaneous-channel formulation address a different operational problem; our finite-budget probe debit and same-stream decision-event correction do not follow from that time-average objective.
Yu et al.~\cite{yu2025sdde}, Sections~III--IV and Propositions~12--14, use an SDDE/arrival-process approximation to optimize worker activation, including a finite-service queue.
That queue model and its delay law are not established by our fixed-delay experiment.

Saad et al.~\cite{saad2023covariance}, Assumption~1, Protocol~1 and Theorems~4.1/5.1, adapt joint arm queries to unknown covariance with independent rounds and a fixed-confidence identification objective.
Our dependence hypotheses are supplied, rather than learned nonparametrically; our comparison instead prices identification through the subsequent learner.
Li et al.~\cite{li2026parameterfree}, Theorem~4.12, remove unknown mixing/eigenvalue inputs from a regularized TD step schedule.
That reduces tuning requirements that our general confidence interface still has, but does not construct a calibrated multi-action terminal-risk shield.
Li et al.~\cite{li2026shrinkingtube}, Assumptions~1--7 and Theorem~1, give shrinking-tube reliability for projected adaptive-kernel SA using uniform drift/minorization and kernel continuity.
Their mixed-history result is broader than a frozen-kernel table; our Gaussian decision correction and paid finite horizon are different outputs, not a replacement for their regularity assumptions.
Reference numbers in this supplement refer to the main paper's bibliography.

\section{Complete Tracking Proof}
Set $R_\star=d/(1-c)$ and $\Gamma_r=(Z_r-R_\star)_+$.
The explicitly assumed uniform-history drift and the positive-part replacement argument give
\begin{align}
\widetilde Z_{r+1}
&\le R_\star+c\Gamma_r,\nonumber\\
\Gamma_{r+1}
&\le(1+\chi)c\Gamma_r+\chi R_\star+(1+\chi^{-1})v^2.
\label{eq:tracking-old-target}
\end{align}
The second line follows by applying Young's inequality to the target shift and subtracting $R_\star$.
Iteration and $\Gamma_0\le Z_0$ prove Eq.~\eqref{eq:tracking-bound} whenever $\bar c=(1+\chi)c<1$.
Finally, $\widehat A_w-A^\star=\{\gamma\phi(s')-\phi(s)\}^\top(w-w_n^\star)$ and $\|\gamma\phi(s')-\phi(s)\|\le(1+\gamma)\Phi$; Cauchy--Schwarz proves Eq.~\eqref{eq:advantage-transfer}.

\section{Complete Occupation-Measure Proof}
Let $\Delta_j(z,a)\ge0$ be the one-step opportunity loss relative to the instance-$j$ certified oracle.
For a stopped predictable policy, define
\begin{equation}
\nu_j(G,a)=\E_j^\pi\sum_{k\le\tau}\mathbf{1}\{Z_k\in G,A_k=a\}.
\end{equation}
Every $\delta$-correct policy induces an occupation measure satisfying the controlled Kalman flow constraints,
\begin{align}
\int\mathcal{I}_{j,1-j}(z,a)\nu_j(dz,da)&\ge\operatorname{kl}(1-\delta,\delta),\nonumber\\
\int(h+q(a))\nu_j(dz,da)&\le B_{\rm msg},\nonumber\\
\int b(a)\nu_j(dz,da)+\ell\Pp_j(\tau>0)&\le B_{\rm env}.
\end{align}
Let $T_a$ be the belief transition and $\zeta_j$ the terminal belief law.
For every bounded measurable $f$, predictability and stopped telescoping give
\begin{align}
\int f(z)\zeta_j(dz)-f(z_1)
&=\int\left[\int f(z')T_a(z,dz')\right.\nonumber\\
&\left.\hspace{18mm}-f(z)\right]\nu_j(dz,da).
\label{eq:belief-flow}
\end{align}
Moreover,
\begin{equation}
\E_j^\pi\sum_{k\le\tau}\Delta_j(Z_k,A_k)
=\int\Delta_j(z,a)\nu_j(dz,da).
\end{equation}
Eq.~\eqref{eq:adaptive-kl}, Eq.~\eqref{eq:budgets}, and Eq.~\eqref{eq:belief-flow} place every admissible occupation measure in the controlled feasible set.
Taking the infimum of its exact loss over that larger set proves Proposition~\ref{prop:occupation}.
\input{transfer_details}

\section{Direct Fixed-Model Validation of the Actual-Risk Gate}
\subsection{Models, learners, and paid protocol}
This study executes the actual-risk gate prospectively on new seeds, rather than reconstructing decisions from the earlier saved study.
The 16 public families cross $Q\in\{8,32\}$, high-instance variance $\theta_1\in\{0.5,2\}$, common-factor persistence $\lambda\in\{0.2,0.8\}$, and withheld-update delay $D\in\{0,8\}$.
The low-instance variance is $\theta_0=0.05$, private variance is one, and message overhead is $h=8$.
Each family is evaluated under both true instances, each learner, and scales $s\in\{32,128,512,2048\}$, with $B_{\rm msg}=(h+Q)s$ and $B_{\rm env}=4s$.
The 128 new seeds are $20571001,\ldots,20571128$, screened against the historical seed registries before execution.
All six policies share policy-independent potential observations within each seed, family, true instance, and learner, including across budget scales.
They do not share potential observations across learners.

The common observation component is stationary Gaussian autoregressive noise, with coefficient $\lambda$ and marginal variance $\theta_j$; each actor has independent unit-variance Gaussian private noise.
An average over the first $q$ actors therefore has private variance $1/q$.
Both learners start from one with a constant-one initial history.
The running mean uses its exact averaging weights; the averaged delayed affine learner uses $\mu=0.5$ and $\eta=0.2/(D+1)$ with the first-in, first-out recurrence analyzed above.
Every scheduled block is paid, including the last $D$ withheld updates: $M$ scheduled updates give $N=\max\{0,M-D\}$ learner updates.
There is no free startup or independent learning reset after a probe.
The measured endpoint is squared parameter error at the exact zero target, not an unobserved within-run convergence curve.

Before data, the fixed probe uses $q_d=Q$, $b_d=1$, and the smallest integer $n_d\in\{1,\ldots,4096\}$ whose exact Gaussian affinity is at most $0.025$.
The log-likelihood threshold is zero, with ties assigned to the low instance.
An infeasible or unavailable probe falls back to the full-budget all-agent program.
After charging the probe, the two commitment programs minimize exact stationary risk under their respective supplied hypotheses over powers-of-two $q\le Q$ and $b\in\{1,2,4,8\}$, breaking ties lexicographically in $(q,b)$.
Both programs are fixed before observing the probe; the true instance is not an input to any deployment decision.
Writing $F_j$ for the full-budget all-agent population risk and $U_j$ for the same-stream population upper bound, the gate requires $U_1<F_1$ and $U_0\le F_0+0.1$.
The no-carryover ablation deletes only the $\sqrt{\delta_j}\Gamma_j$ contribution from that admission calculation.
The other controls are information-only probing, all-agent fallback, fixed $(q,b)=(2,1)$, and a known-instance oracle used only for evaluation.

The complete grid contains $16\times2\times2\times4\times6\times128=196608$ trajectories.
Every trajectory is regenerated in a full semantic replay, including its paid ledgers and decisions.
Two analyses reproduce all seven output tables byte for byte.
No failed, infeasible, adverse, or partial-cost record is removed.

\subsection{All macro comparisons and counterevidence}
For each learner, instance, and scale, we average the paired fallback-minus-policy squared-risk difference across all 16 families within each seed.
Table~\ref{tab:direct-macro} reports all 16 corrected-gate macro means and descriptive percentile intervals from 2,000 paired-seed bootstrap resamples.
These are cellwise descriptive intervals, not a familywise significance or equivalence claim.
\begin{table*}[t]
\centering
\caption{All corrected-gate macro risk improvements over all-agent fallback.
Positive values favor the gate; brackets give descriptive 95 percent paired-bootstrap intervals.}
\label{tab:direct-macro}
\begin{tabular}{llrr}
\toprule
Learner & Scale & High instance: mean [interval] & Low instance: mean [interval] \\
\midrule
Running mean & 32 & $0.07820\ [0.06227,0.09303]$ & $-0.002290\ [-0.003226,-0.001357]$ \\
& 128 & $0.03452\ [0.03051,0.03840]$ & $0.000763\ [0.000530,0.000986]$ \\
& 512 & $0.009436\ [0.008400,0.010491]$ & $0.000373\ [0.000324,0.000422]$ \\
& 2048 & $0.002268\ [0.001998,0.002550]$ & $0.0000860\ [0.0000736,0.0000985]$ \\
Delayed affine & 32 & $0.10001\ [0.08038,0.11922]$ & $0.06114\ [0.05777,0.06469]$ \\
& 128 & $0.14696\ [0.13225,0.16261]$ & $0.07750\ [0.07538,0.07955]$ \\
& 512 & $0.04006\ [0.03541,0.04538]$ & $0.010206\ [0.009663,0.010712]$ \\
& 2048 & $0.008698\ [0.007723,0.009715]$ & $0.000914\ [0.000838,0.000993]$ \\
\bottomrule
\end{tabular}
\end{table*}

All eight high-instance macro intervals are positive, as are seven of eight low-instance intervals.
The exception is the low-instance running mean at scale 32, where probing can spend resources without compensating risk reduction.
At the finer family level, 23 of 256 corrected-gate mean improvements are negative (four high and nineteen low); 174 intervals are strictly positive, nine strictly negative, and 73 include zero or are exactly zero.
The permitted low-instance population-risk increase is $0.1$, so this adverse result does not contradict a zero-harm condition: no such condition was imposed.

The corrected gate admits probing on 113 of 128 public designs, exactly as does the no-carryover gate.
Their probe-use indicators, actions, and risks coincide on all 32,768 paired policy rows.
The selection-dependence term remains necessary to the stated bound, but its incremental empirical decision benefit is unresolved on this grid.
Small differences between separately resampled bootstrap endpoints for the two identical policies do not constitute performance differences.
The prespecified zero-allowance admission diagnostic admits 100 of 128 designs; it uses saved branch calculations, not additional trajectories.

Table~\ref{tab:direct-costs} retains all six policies' pooled risks and paid costs.
Pooling gives equal weight to heterogeneous families, instances, learners, and scales, so it is descriptive rather than a universal ranking or a comparison at identical realized costs.
Fixed $(q,b)=(2,1)$ has lower pooled risk and environment cost than the corrected gate, whereas the gate sends fewer messages.
The oracle is an instance-aware reference and is not deployable without the unknown label.
\begin{table*}[t]
\centering
\caption{Descriptive pooled means over all 256 cells for each policy; every cell has 128 seeds.}
\label{tab:direct-costs}
\begin{tabular}{lrrrr}
\toprule
Policy & Squared risk & Paid messages & Paid environment & Probe fraction \\
\midrule
All agents & $0.114864$ & $19040.000$ & $680.000$ & $0$ \\
Fixed $(2,1)$ & $0.062931$ & $19037.500$ & $1903.750$ & $0$ \\
Information only & $0.099064$ & $18509.108$ & $2467.163$ & $0.937500$ \\
Oracle (evaluation only) & $0.051335$ & $18508.477$ & $2509.625$ & $0$ \\
Without carryover & $0.079311$ & $18509.341$ & $2463.203$ & $0.882813$ \\
Corrected gate & $0.079311$ & $18509.341$ & $2463.203$ & $0.882813$ \\
\bottomrule
\end{tabular}
\end{table*}

\subsection{Probe validity and the scope of the enclosures}
There are 8,192 distinct deduplicated potential probe decisions and thirteen observed classification errors, an empirical fraction of $0.001587$.
The exact affinity bounds range from $0.021139$ to $0.024380$, each below $0.025$ under both supplied hypotheses.
This new fixed-probe construction does not retroactively validate the earlier probe's nominal $0.025$ directional interpretation, which bounded equal-prior average error instead.
The historical reanalysis continues to use its larger exact affinities.

The risk enclosures concern population expectations under known binary Gaussian hypotheses and deterministic linear programs.
In 123 of 256 corrected-gate cells, the finite sample mean exceeds the corresponding population upper bound; such an exceedance is not by itself a population-bound violation.
In particular, fallback risk is an exact model expectation, not a simultaneous upper confidence limit for noisy sample means.
The study therefore does not establish samplewise coverage, calibration for unknown kernels, or a general neural policy-learning guarantee.

Figs.~\ref{fig:direct-rm-low}--\ref{fig:direct-affine-high} show every policy against budget scale, actual paid messages, and actual paid environment units, separately for each learner and instance.
Each point averages all 16 families and all 128 seeds; connecting points does not imply a within-trajectory learning curve or equal realized cost.
\begin{figure*}[t]
\centering
\includegraphics[width=\textwidth]{figures/direct_gate_running_mean_low_20261020.pdf}
\caption{Direct validation, running mean and low instance: all-policy risk versus allocated scale and actual paid resources. Corrected and no-carryover curves coincide.}
\label{fig:direct-rm-low}
\end{figure*}
\begin{figure*}[t]
\centering
\includegraphics[width=\textwidth]{figures/direct_gate_running_mean_high_20261020.pdf}
\caption{Direct validation, running mean and high instance, with the same complete-family averaging and all six policies.}
\label{fig:direct-rm-high}
\end{figure*}
\begin{figure*}[t]
\centering
\includegraphics[width=\textwidth]{figures/direct_gate_averaged_delayed_affine_low_20261020.pdf}
\caption{Direct validation, averaged delayed affine learner and low instance; every scheduled update is charged, including withheld terminal updates.}
\label{fig:direct-affine-low}
\end{figure*}
\begin{figure*}[t]
\centering
\includegraphics[width=\textwidth]{figures/direct_gate_averaged_delayed_affine_high_20261020.pdf}
\caption{Direct validation, averaged delayed affine learner and high instance; exact overlap of the two gates is retained without graphical offsets.}
\label{fig:direct-affine-high}
\end{figure*}
\end{document}
